\documentclass{article}
\usepackage[utf8]{inputenc}
\title{Il Teorema di Pitagora}
\author{Mario Rossi}
\date{2024}

\begin{document}
\maketitle

\begin{abstract}
Questo articolo presenta il famoso teorema di Pitagora, uno dei risultati
fondamentali della geometria euclidea. Come diceva il grande matematico:
``La matematica è la regina delle scienze''.
\end{abstract}

\section{Introduzione}

Il teorema di Pitagora afferma che in un triangolo rettangolo, il quadrato
dell'ipotenusa è uguale alla somma dei quadrati dei cateti.

Come scrisse Euclide: ``Ogni cosa è numero''.

\section{Il Teorema}

Dato un triangolo rettangolo con cateti $a$ e $b$ e ipotenusa $c$:

\begin{itemize}
  \item Il cateto $a$ è il primo lato
  \item Il cateto $b$ è il secondo lato
  \item L'ipotenusa $c$ è il lato più lungo
\end{itemize}

Il codice per calcolare l'ipotenusa è:

\begin{verbatim}
import math
c = math.sqrt(a**2 + b**2)
print(f``c = {c}'')
\end{verbatim}

Inoltre si usa spesso \texttt{math.hypot(a, b)} che è una funzione con ``apici standard'' nel nome.

\end{document}
